% voorbeelddocument voor Inleiding Informatiekunde
% alles wat volgt achter een procentteken is commentaar 
% en wordt genegeerd bij de verwerking van het document

% definitie van de documentklasse (artikel) en de fontgrootte (11 punten) 
\documentclass[11pt]{article}
\usepackage[utf8]{inputenc}  % gebruik de juiste 'character encoding'
\usepackage[dutch]{babel}    % definitie van de taal (Engels is de standaard)
\usepackage{hyperref}        % geef URLs netjes weer
\usepackage{booktabs}		 % mooiere tabellen
\usepackage{a4wide}          % papierformaat en marges
\pagestyle{empty}            % zet paginanummering uit

% titel en auteur van het document
\title{Titel van dit verslag}
\author{Naam van de auteur}
% \date{} % de datum van het document 

% einde definities; start van de tekst
\begin{document}

\maketitle                  % maakt een kopje met titel, auteur en datum
\thispagestyle{empty}       % zet paginanummering uit voor deze pagina

% kopje met bijbehorende titel; varianten: \subsection, \subsubsection
\section*{Inleiding}

% tekst
De inleiding bevat uitleg over wat online adverteren is, en wat de voor- en nadelen zijn van online advertenties en aanbevelingen. Je inleiding bevat in ieder geval een vraagstelling. Leg verder uit waarom je denkt dat het belangrijk of interessant is om antwoord te krijgen op je onderzoeksvraag. Zie de college-aantekeningen voor meer informatie over wat je in de inleiding kunt vermelden, en/of Nederhoed, hoofdstuk 10, en/of \url{http://practicumav.nl/onderzoeken/onderzoeksverslag.html}

Wanneer je tekst van anderen wilt gebruiken, moet je dit als een citaat markeren, bevoorbeeld als volgt: 

\begin{quote}
The basic principle of recommendations is that significant dependencies exist between userand
item-centric activity. (Aggarwal, 2016)
\end{quote}


\section*{Methode}

Leg uit van welke data je gebruik hebt gemaakt, en hoe je die data gemaakt of verzameld hebt. Leg daarna uit hoe je de data geanalyseerd hebt. Wat heb je bijvoorbeeld uitgerekend, en wat heb je met elkaar vergeleken.

\section*{Resultaten }

Geef aan wat de uitkomsten van je onderzoek zijn.

\section*{Discussie}

Bespreek hier wat het antwoord op je onderzoeksvraag is, en of dit iets betekent voor verder onderzoek naar bijvoorbeeld het nut van online adverteren. Je kunt ook iets zeggen over de tekortkomingen van je eigen onderzoek (wat had er beter gekund?) en mogelijkheden tot vervolgonderzoek. 

Je verslag is  \textbf{500 \`a 600 woorden} lang,  de lijst met bronnen niet meegerekend. 


ewpage
\section*{Bronnen}

Geef een lijst met alle gebruikte bronnen.

\begin{itemize}
\item C. Aggarwal, An Introduction to Recommendere Systems, Springer, 2016
\item \url{http://www.youronlinechoices.com/nl/interest-based-advertising}
\end{itemize}

\end{document}
